% day3_seg2_inertial/impinst.tex
\begin{frame}
  \onslide<1->{
	\frametitle{Inertial Error Dynamics} 
	\framesubtitle{Imperfect Inertial Instruments}
}
\begin{itemize}
  \onslide<2->{
 \item Accelerometers and Gyros have inherent imperfections
}
  \onslide<3->{
 \item Bias, scale factor, $g^2$ terms, noise, random walks, etc.
}
  \onslide<4->{
 \item Inertial sensor clusters also have errors: misalignments
}
  \onslide<5->{
 \item Many error sources are calibrated out using rate tables during ATP
}
 \onslide<6->{ 
 \item Residual effects will influence strapdown nav equations
}
\end{itemize}
\end{frame}

\begin{frame}
  \onslide<1->{
	\frametitle{Inertial Error Dynamics} 
	\framesubtitle{Imperfect Inertial Instruments}
}
\begin{itemize}
  \onslide<2->{
  \item Strapdown nav equations very sensitive to gyro errors
}
  \onslide<3->{
  \item Gyro biases yield erroneous $\dcm{B}{L}$ matrix which in turn causes erroneous accelerations to be computed in local level frame
}
  \onslide<4->{
  \item Thus, gyro errors affect not only attitude, but also position and velocity
}
\end{itemize}
\end{frame}

\begin{frame}
  \onslide<1->{
	\frametitle{Inertial Error Dynamics} 
	\framesubtitle{Simple Inertial Error Models}
}
\begin{itemize}
  \onslide<2->{
  \item Gyros
\[
\simprateerreq{x}{y}{z}
\]
\[
\simprateerreq{y}{z}{x}
\]
\[
\simprateerreq{z}{x}{y}
\]

where $\gyrobias{i}$ is the $i^\text{th}$ gyro bias, $\gyrosf{i}$ the $i^\text{th}$ gyro scale factor, and $\gyromisalign{i}{j}$ is the misalignment of the $i^\text{th}$ gyro axis toward the $j^\text{th}$ gyro axis. 
}
\end{itemize}
\end{frame}

\begin{frame}
  \onslide<1->{
	\frametitle{Inertial Error Dynamics} 
	\framesubtitle{Simple Inertial Error Models}
}
\begin{itemize}
  \onslide<2->{
  \item Accelerometers
\[
\simpaccelerreq{x}{y}{z}
\]
\[
\simpaccelerreq{y}{z}{x}
\]
\[
\simpaccelerreq{z}{x}{y}
\]
where $\accelbias{i}$ is the $i^\text{th}$ accelerometer bias, $\accelsf{i}$ the $i^\text{th}$ accelerometer scale factor, and $\accelmisalign{i}{j}$ is the misalignment of the $i^\text{th}$ accelerometer axis toward the $j^\text{th}$ accelerometer axis. 
}
\end{itemize}
\end{frame}


\begin{frame}
  \onslide<1->{
	\frametitle{Inertial Error Dynamics} 
	\framesubtitle{Bias Models}
}
\begin{itemize}
  \onslide<2->{
  \item Bias in both types of instruments can be further broken down into several components:
  \begin{itemize}
	\item fixed bias -- constant from turn-on to turn-on -- should be calibrated out
	\item bias stability -- varies from turn-on to turn-on -- can't be calibrated out, but can be characterized in terms of variance -- can model as random constant
	\item bias drift -- usually modeled as random walk
  \end{itemize} 
}
\end{itemize}
\end{frame}


\begin{frame}
  \onslide<1->{
	\frametitle{Inertial Error Dynamics} 
	\framesubtitle{Scale Factor Models}
}
\begin{itemize}
  \onslide<2->{
  \item Scale factor can also be modeled as a sum of three components
}
  \onslide<3->{
  \item Often in nav filters, scale factor residual is modeled as a random constant
}
  \onslide<4->{
  \item Important to keep in mind that the contribution of scale factor error is a function of input (rate or acceleration) 
}
\end{itemize}
\end{frame}
\begin{frame}
  \onslide<1->{
	\frametitle{Inertial Error Dynamics} 
	\framesubtitle{Higher Order Models}
}
\begin{itemize}
  \onslide<2->{
  \item In expanding the output error as a Taylor series function of input, bias, scale factor, and misalignment show up as zero and first order terms.
}
  \onslide<3->{
  \item Higher order terms are also present, especially second order terms.
}
  \onslide<4->{
  \item In some calibration and alignment filters, such higher order effects are modelled
}
  \onslide<5->{
  \item SNU-84-1 includes error budgets associated with all the combined, higher order effects.  
}
\end{itemize}
\end{frame}

\begin{frame}
  \onslide<1->{
	\frametitle{Inertial Error Dynamics} 
	\framesubtitle{Higher Order Models:  Taylor Series of Two Variables}
}
  \onslide<2->{
\begin{align*}
\gentaylordouble{f}{x}{y}{x_0}{y_0}
\end{align*}
Taylor Series provides basis for all error analysis.
}
\end{frame}
